\documentclass[a4paper,11pt]{article}
\usepackage[utf8]{inputenc}
\usepackage[T1]{fontenc}
\usepackage[english]{babel}
\usepackage[left=2.5cm,right=2.5cm,top=2cm,bottom=2cm]{geometry}
\usepackage{hyperref}
\usepackage{xcolor}
\usepackage{microtype}
\usepackage{lmodern}

\definecolor{primary}{RGB}{0, 51, 102}
\hypersetup{colorlinks=true, urlcolor=primary}

\begin{document}
\noindent
\textbf{Loïc Aron MBASSI EWOLO} \\
ENSEA -- Double Diplôme Ingénieur \\
\href{mailto:loicaron.mbassiewolo@ensea.fr}{loicaron.mbassiewolo@ensea.fr} \quad (+33) 7 59 52 33 98 \\
\href{https://github.com/Nameless0l}{github.com/Nameless0l} \quad \href{https://mbassiloic.tech}{mbassiloic.tech}

\vspace{1em}
\noindent Cergy, le 22 mars 2026

\vspace{1em}
\noindent
\textbf{Objet : Candidature -- Data Scientist} \\
\textbf{Référence : Offre d'alternance -- Framatome}

\vspace{1.5em}
\noindent Madame, Monsieur,

\vspace{0.8em}
Le traitement intelligent de données massives pour l'optimisation des processus critiques est au cœur de ma formation d'ingénieur à ENSEA et de mon stage en cours à INRIA Saclay. Vous recherchez un Data Scientist capable de transformer les données du secteur nucléaire en intelligence opérationnelle. Ma spécialisation en Machine Learning appliqué, combinée à une expérience concrète de pipeline NLP en environnement de recherche d'excellence, me positionne pour contribuer rapidement à vos projets data.

\vspace{0.8em}
Au sein de l'équipe TRIBE à INRIA Saclay (avril-août 2026), j'ai développé un pipeline complet d'analyse de données massives en Python, intégrant NLP, traitement de données structurées et évaluation statistique. Ce travail m'a permis de maîtriser pandas, scikit-learn et les bonnes pratiques de validation de modèles. Auparavant, au programme UROP 2024 (reconnaissance des anomalies cardiaques via électrocardiogrammes), j'ai conçu et entraîné un modèle de classification basé sur TensorFlow/Keras, incluant extraction de features, optimisation d'hyperparamètres et gestion d'imbalance de données. Ces deux expériences m'ont inculqué la rigueur scientifique indispensable à la data science critique.

\vspace{0.8em}
Ma participation au premier prix du hackathon national CITS (Urban Waste Management) démontre ma capacité à appliquer le Machine Learning à des problématiques réelles : prédiction des besoins de collecte, optimisation des routes via algorithmes combinatoires et analyse coût-bénéfice. J'ai manipulé des données hétérogènes (IoT, météo, démographiques) et proposé des solutions quantifiables. Ces projets reflètent mon approche : au-delà de l'algorithme, j'accorde de l'importance au contexte métier et à la mesurabilité des résultats.

\vspace{0.8em}
Je suis disponible dès septembre 2026 pour une alternance de 12 mois. Passionné par les défis data du secteur critique, je suis convaincu que mon expérience INRIA, ma formation ENSEA et mon implication dans des projets concrets feront de moi un atout pour Framatome dans la valorisation de ses données.

\vspace{1.5em}
\noindent Veuillez agréer, Madame, Monsieur, l'expression de mes salutations distinguées.

\vspace{2em}
\noindent \textbf{Loïc Aron MBASSI EWOLO}

\end{document}
